% =============================================================================
% Chapter 4: How to Use the IP
% =============================================================================
\chapter{How to Use the IP}
\label{ch:howto}

This chapter provides step-by-step guidance for software developers and firmware engineers who need
to configure and operate the APB SPI Master IP from a processor.

\section{Basic Configuration Sequence}

Before starting any SPI transaction, the following registers must be configured:

\begin{enumerate}[leftmargin=2em]
  \item \textbf{Set the SPI clock divider} (CLKDIV, offset 0x04):
  \begin{lstlisting}
  apb_write(REG_CLKDIV, 8'd4);
  // SPI clock period = 2 * (4+1) * 10ns = 100ns (10 MHz at 100 MHz sys clk)
  \end{lstlisting}

  \item \textbf{Set the command value and length} (SPICMD at 0x08, SPILEN bits [5:0]):
  \begin{lstlisting}
  apb_write(REG_SPICMD, 32'h0000_0003);  // e.g., READ command
  // cmd_len = 8 bits (set in SPILEN[5:0])
  \end{lstlisting}

  \item \textbf{Set the address value and length} (SPIADR at 0x0C, SPILEN bits [13:8]):
  \begin{lstlisting}
  apb_write(REG_SPIADR, 32'h0000_00AB);  // Target address
  // addr_len = 8 bits (set in SPILEN[13:8])
  \end{lstlisting}

  \item \textbf{Set the data length} (SPILEN bits [31:16]):
  \begin{lstlisting}
  apb_write(REG_SPILEN, {16'd8, 2'b00, 6'd8, 2'b00, 6'd8});
  // data_len=8, addr_len=8, cmd_len=8 (all in bits)
  \end{lstlisting}

  \item \textbf{Set dummy cycles if needed} (SPIDUM at 0x14):
  \begin{lstlisting}
  apb_write(REG_SPIDUM, {16'd0, 16'd4});
  // dummy_wr=0, dummy_rd=4 cycles
  \end{lstlisting}
\end{enumerate}

\section{Performing a Standard SPI Write}

\begin{enumerate}[leftmargin=2em]
  \item Configure registers as above with \texttt{cmd\_len=0}, \texttt{addr\_len=0} for a simple data-only write.
  \item Load data into the TX FIFO:
  \begin{lstlisting}
  apb_write(REG_TXFIFO, 32'hA5A5_A5A5);
  \end{lstlisting}
  \item Trigger the write operation:
  \begin{lstlisting}
  apb_write(REG_STATUS, {20'h0, 4'b0001, 4'h0, 4'b0010});
  // cs=CS0 (bit 8), op=wr (bit 1)
  \end{lstlisting}
  \item Wait for EOT by polling \texttt{events\_o[1]} or using an interrupt.
\end{enumerate}

\section{Performing a Standard SPI Read}

\begin{enumerate}[leftmargin=2em]
  \item Configure registers with desired data length.
  \item Trigger the read operation:
  \begin{lstlisting}
  apb_write(REG_STATUS, {20'h0, 4'b0001, 4'h0, 4'b0001});
  // cs=CS0, op=rd (bit 0)
  \end{lstlisting}
  \item Wait for EOT.
  \item Read data from the RX FIFO:
  \begin{lstlisting}
  apb_read(REG_RXFIFO, rx_data);  // Pops one 32-bit word
  \end{lstlisting}
\end{enumerate}

\textbf{Important:} In standard SPI mode, the IP samples \texttt{sdi1} (not \texttt{sdi0}) for
receive data. Ensure the slave device drives the correct SDI line.

\section{Performing a Quad SPI Write}

Same as standard write, but use \texttt{op = 4'b1000} (qwr bit [3]). The IP will drive all four
\texttt{sdo0--sdo3} lines simultaneously, achieving 4$\times$ throughput.

\section{Performing a Quad SPI Read}

Same as standard read, but use \texttt{op = 4'b0100} (qrd bit [2]). The IP samples all four
\texttt{sdi0--sdi3} lines simultaneously.

\section{CMD + ADDR + DUMMY + DATA Transactions}

For flash memory-style transactions, configure non-zero lengths for command, address, and dummy phases:
\begin{lstlisting}
// Read 8 bits with CMD=0x03, ADDR=0xAB, 4 dummy cycles
cfg_spi(8'd4,
    32'h0000_0003, 6'd8,    // cmd=0x03, 8 bits
    32'h0000_00AB, 6'd8,    // addr=0xAB, 8 bits
    16'd8,                   // data_len=8 bits
    16'd4, 16'd0);           // dummy_rd=4, dummy_wr=0
trigger(4'b0001);            // Standard read
\end{lstlisting}

The FSM will automatically sequence: CMD $\rightarrow$ ADDR $\rightarrow$ DUMMY $\rightarrow$ DATA\_RX.

\section{Multi-Word Transfers}

For transfers longer than 32 bits, set \texttt{data\_len} to the total number of bits. The FIFO
automatically provides/consumes 32-bit words as needed:
\begin{lstlisting}
// TX burst: 4 words x 32 bits = 128 bits
cfg_spi(8'd2, 32'h0, 6'd0, 32'h0, 6'd0, 16'd32, 16'd0, 16'd0);
push_tx(32'hDEAD_BEEF);
push_tx(32'hCAFE_BABE);
push_tx(32'h1234_5678);
push_tx(32'hABCD_EF01);
// Trigger 4 separate 32-bit writes, or one 128-bit write
\end{lstlisting}

For multi-word RX, read multiple words from RXFIFO after EOT:
\begin{lstlisting}
cfg_spi(8'd2, 32'h0, 6'd0, 32'h0, 6'd0, 16'd64, 16'd0, 16'd0);
trigger(4'b0001);  // Read 64 bits
wait_eot();
pop_rx(word0);     // First 32 bits
pop_rx(word1);     // Next 32 bits
\end{lstlisting}

\section{Software Reset}

To clear both FIFOs and abort any pending operation:
\begin{lstlisting}
apb_write(REG_STATUS, 32'h0000_0010);  // Set bit [4] = swrst
\end{lstlisting}

After the software reset, the STATUS register should read back with the idle bit set and both
FIFO element counts at zero. Note that the software reset clears FIFOs via \texttt{clr\_i} but
does \emph{not} affect configuration registers (CLKDIV, SPICMD, etc.).

\textbf{Caveat:} The controller FSM's \texttt{spi\_swrst} input has a FIXME comment in the
RTL---it is not used within the controller itself. The reset only clears the FIFOs at the
top level. If the controller is mid-transaction, the behavior is undefined. Best practice:
wait for EOT before issuing software reset.

\section{Interrupt Configuration}

To enable interrupts:
\begin{lstlisting}
apb_write(REG_INTCFG, 32'h8000_0000);  // Set bit [31] = int_en
\end{lstlisting}

The interrupt configuration register fields:
\begin{table}[H]
\centering
\caption{INTCFG Register Fields}
\begin{tabularx}{\textwidth}{@{}clX@{}}
\toprule
\textbf{Bits} & \textbf{Field} & \textbf{Description} \\
\midrule
{[LOG\_BUFFER\_DEPTH:0]}     & int\_th\_tx  & TX FIFO threshold (interrupt when elements $\leq$ threshold). \\
{[8+LOG\_BUFFER\_DEPTH:8]}   & int\_th\_rx  & RX FIFO threshold (interrupt when elements $\geq$ threshold). \\
{[16+LOG\_BUFFER\_DEPTH:16]} & int\_cnt\_tx & TX counter target for periodic interrupts. \\
{[24+LOG\_BUFFER\_DEPTH:24]} & int\_cnt\_rx & RX counter target for periodic interrupts. \\
{[30]}                       & int\_cnt\_en & Enable counter-based interrupt mode. \\
{[31]}                       & int\_en      & Global interrupt enable. \\
\bottomrule
\end{tabularx}
\end{table}

\textbf{Safe interrupt enable sequence:}
\begin{enumerate}[leftmargin=2em]
  \item Assert system reset.
  \item Release reset, wait for stabilization.
  \item Optionally read \texttt{REG\_INTSTA} to arm the interrupt FSM.
  \item Write \texttt{REG\_INTCFG} with desired thresholds and \texttt{int\_en = 1}.
\end{enumerate}

\section{Chip-Select Selection}

To address a specific slave device, set the appropriate bit in \texttt{spi\_csreg} (STATUS[11:8]):
\begin{lstlisting}
// Select CS1 for a write operation
apb_write(REG_STATUS, {20'h0, 4'b0010, 4'h0, 4'b0010});
// csreg=0010 -> CSN1 active, CSN0/2/3 stay high
\end{lstlisting}

Multiple CS bits can be set simultaneously to address multiple slaves (broadcast), though this
is rarely used in practice.
